\chapter{Knocks and Bumps}

\section*{Definition and Overview}
Knocks, bumps, bruises, and minor lumps are everyday childhood and adult experiences that occur when the body is struck against something or something strikes the body. A knock to the head, limb, or trunk can cause bruising (a collection of blood under the skin), swelling, and pain. Most knocks are minor and settle quickly with simple first aid, but some can be serious, particularly knocks to the head, the abdomen, or the eye, or those that cause significant deformity or loss of function. This chapter explains how to assess a knock, treat it at home, and recognise when medical attention is needed.

\section*{Epidemiology}
Knocks and bumps are among the most common injuries at all ages. In children they are a normal part of exploration, play, and sport, and most cause nothing more than a bruise. Falls are the leading cause of injury presentations in young children and older adults, and head knocks are particularly common in both groups. Sporting injuries, workplace accidents, and household bumps account for most of the remainder. Although the vast majority are minor, a small proportion cause fractures, internal injuries, or concussion, so knowing how to tell the difference matters.

\section*{Aetiology and Pathophysiology}
\begin{itemize}
    \item \textbf{Blunt trauma:} a blow from a fall, collision, or falling object damages small blood vessels beneath the skin
    \item \textbf{Bruising:} blood leaks from damaged capillaries into the tissues, producing the familiar black-and-blue mark that changes colour over days
    \item \textbf{Swelling:} inflammation and fluid accumulation in the injured area
    \item \textbf{Contusions:} bruising of deeper structures, including muscle, bone, or internal organs
    \item \textbf{Haematoma:} a larger collection of blood that forms a lump, such as a "goose egg" on the forehead
    \item \textbf{Head knocks:} may cause concussion, which is a temporary disturbance of brain function
    \item \textbf{Bone injury:} a hard knock may cause a fracture even without obvious deformity
    \item \textbf{Risk factors:} anticoagulant medication, bleeding disorders, and older age increase the tendency to bruise
\end{itemize}

\section*{Clinical Features}
\begin{itemize}
    \item Pain and tenderness at the site of the knock
    \item Bruising, which may appear immediately or over the following day
    \item Swelling, sometimes with a visible lump
    \item Redness or warmth over the area
    \item Reduced movement of an affected limb
    \item With head knocks: headache, dizziness, or confusion
    \item With abdominal knocks: pain, tenderness, or vomiting
\end{itemize}

\section*{Differential Diagnosis}
\begin{itemize}
    \item \textbf{Fracture:} suspected when there is severe pain, deformity, inability to bear weight, or tenderness over a bone
    \item \textbf{Concussion:} after a head knock with confusion, vomiting, or loss of consciousness
    \item \textbf{Internal injury:} after a significant blow to the abdomen or chest
    \item \textbf{Sprains and strains:} injury to ligaments or muscles near the site of the knock
    \item \textbf{Eye injury:} a knock to the eye warrants assessment for damage to the eyeball or surrounding bone
    \item \textbf{Unexplained bruising:} may indicate a bleeding disorder, medication effect, or, in children, possible abuse
\end{itemize}

\section*{When to Seek Medical Care}
Most knocks can be managed at home, but see a doctor if there is severe or worsening pain, inability to use the injured part, a large or rapidly expanding bruise, bruising that appears without a clear cause, or symptoms that do not settle within a few days.

\textbf{Call 000 (or local emergency services) for:}
\begin{itemize}
    \item A head knock with loss of consciousness, confusion, vomiting, seizures, or unequal pupils
    \item A blow to the abdomen with severe pain, tenderness, or vomiting
    \item An obviously deformed or angulated limb or inability to move or bear weight on it
    \item A knock that penetrates the skin or causes heavy bleeding
    \item A knock to the eye with severe pain, vision loss, or blood within the eye
\end{itemize}

\section*{Diagnosis and Investigations}
\begin{itemize}
    \item \textbf{History:} how the knock happened, the force involved, and whether there was loss of consciousness or a period of confusion
    \item \textbf{Examination:} assessment of the injury site, range of movement, and neurovascular status
    \item \textbf{X-ray:} when a fracture is suspected
    \item \textbf{CT scan:} after significant head knocks with warning signs, or abdominal knocks where internal injury is possible
    \item \textbf{Observation:} children and adults with head knocks are often monitored at home or in hospital for a period
    \item \textbf{Blood tests:} occasionally used to assess bleeding disorders or anticoagulation
\end{itemize}

\section*{Management}
\begin{itemize}
    \item \textbf{Rest:} rest the injured area and avoid activity that aggravates pain
    \item \textbf{Ice:} apply a cold pack wrapped in a cloth for 15--20 minutes at a time during the first day to reduce swelling
    \item \textbf{Compression:} a firm bandage may help limit swelling
    \item \textbf{Elevation:} raise the injured limb to reduce swelling
    \item \textbf{Pain relief:} paracetamol or ibuprofen as needed
    \item \textbf{Bruise care:} bruises resolve over one to two weeks, changing colour as the blood is reabsorbed
    \item \textbf{Observation after head knocks:} monitor for warning signs for 24 hours and avoid sport until cleared
    \item \textbf{Medical review:} for suspected fracture, concussion, or injuries that fail to settle
\end{itemize}

\section*{Complications}
\begin{itemize}
    \item Compartment syndrome, a rare but serious complication of limb injury causing severe pain, swelling, and impaired blood flow
    \item Undetected fracture or internal injury if a knock is dismissed
    \item Concussion with persisting symptoms, including headache, poor concentration, and mood changes
    \item Large haematomas that may require drainage
    \item Infection if the skin is broken
    \item Frequent unexplained bruising may point to an underlying bleeding disorder
\end{itemize}

\section*{Prognosis and Outcomes}
The vast majority of knocks and bumps heal completely within days to a couple of weeks. Bruises and soft tissue injuries follow a predictable pattern of resolution, and simple first aid speeds recovery. Injuries that cause fractures, concussion, or internal damage require specific treatment, but with appropriate care outcomes are generally excellent. The key is recognising which knocks need more than home care.

\section*{Prevention and Self-Care}
\begin{itemize}
    \item Child-proof the home: soft corner guards, safety gates, and secure furniture reduce knocks in young children
    \item Supervise young children during play and on playground equipment
    \item Use appropriate safety gear for sport, including helmets for cycling, skating, and contact sports
    \item Prevent falls in older adults with handrails, good lighting, non-slip surfaces, and balance exercises
    \item Use ladders safely and secure rugs and cords
    \item Apply first aid promptly after a knock to limit swelling
    \item Seek medical advice if a knock seems more serious than a routine bump
\end{itemize}

\section*{Sources and Further Reading}
The following reputable sources provide further information for healthcare students and patients seeking reliable, current advice about knocks, bumps, and bruising.
\begin{itemize}
    \item Healthdirect. \textit{Bruises and bumps.} https://www.healthdirect.gov.au/bruises
    \item Kids Health (Australia). \textit{Head knocks in children.} https://www.kidshealth.org.nz/
    \item St John Ambulance Australia. \textit{First aid for bumps and bruises.} https://stjohn.org.au/
    \item NHS. \textit{Bumps and bruises.} https://www.nhs.uk/conditions/bumps-and-bruises/
    \item Better Health Channel (Victoria). \textit{Bruising.} https://www.betterhealth.vic.gov.au/
    \item MSD Manual. \textit{Bruising.} https://www.msdmanuals.com/
\end{itemize}
